\chapter{Conclusion}
\label{ch:conclusion}

This thesis asked what happens when language models of different capability share an economy, and
returned an answer that inverts the obvious one. The capable models do not reliably dominate the
weak. With no channel between them the strong earn a premium; but when each agent is asked to make
a public pledge before it acts, the premium turns negative and the strong earn less. The pledge is
never delivered to anyone. What changes behavior is the act of committing in public, and only the
capable models treat their own stated intention as binding: they cooperate on it, contribute
visibly, and are free-ridden by the others through the shared pool. The effect belongs to the group
game and fades in the two-player exchange, which places it in the structure of the commons rather
than in capability itself.

The second finding is a measure rather than a phenomenon. Asked whether the agents deceive, the
obvious divergence between stated and revealed policy gives the wrong answer, calling the steadiest
cooperators the least honest; a contemporaneous, action-by-action surprisal gives the right one. The
contribution is in reporting the two together, the signal and the artefact that mimics it, rather
than quietly keeping the measure that flatters the story. Both results were fixed in direction
before the first game ran, which is the strongest thing the thesis can say for them: they are not
what it set out to find. For anyone assembling systems out of agents of unequal capability, that is
the warning worth keeping: the strong do not simply win, and the familiar ways of managing them,
asking them to be transparent and checking whether they lie, can misfire unless the structure they
act in and the measure used on them are chosen with care.
